\scanpage{226}
\mspage{81}

n'est pas déduite d'une relation entre flèches de $M$ i.e d'une partie de $\mathrm{Fl}(M)\times\mathrm{Fl}(M)$, mais est de nature plus complexe. Pour la décrire, il faut introduire l'ensemble $D$ des diagrammes \inserted{carrés comm. (traits pleins)} de la forme

\diagram{\begin{tikzcd}[ampersand replacement=\&]
{}\arrow[d,phantom,"{(\unsure{\alpha})}" description] \& A\arrow[r]\arrow[d,"i"] \& X\arrow[d,"p"]\\
 {} \& B\arrow[r]\arrow[ur,dashed,"{\unsure{f|g}}"] \& Y
\end{tikzcd}}

(qu'on peut interpréter comme $\mathrm{Fl}(\underline{\mathrm{Fl}}(M))$, ensemble des flèches dans $\underline{\mathrm{Fl}}(M)$)\,;

On a une partie $D_0$ de $D$, \unsure{formée} des diagrammes pour lesquels existe $f$ avec deux triangles commutatifs (c'est donc l'image de $\underset{\mathrm{Hom}(\unsure{\mathbb{\Delta}|\Delta}^{3},M)}{\underline{\mathrm{Fl}_3(M)}}$ dans $\mathrm{Hom}(\mathbb{\Delta}^1\times\mathbb{\Delta}^1,M)$, transposée de $\mathbb{\Delta}^1\times\mathbb{\Delta}^1\longrightarrow\mathbb{\Delta}^3$ donné par $(0,0)\mapsto0$, $(0,1)\mapsto1$, $(1,0)\unsure{\vdash|\mapsto}2$, $(1,1)\vdash3$), et une application

\diagram{\begin{tikzcd}[ampersand replacement=\&,row sep=1.1em]
 D\arrow[r,"{\unsure{\varphi|\phi}}"] \& \mathrm{Fl}(M)\times\mathrm{Fl}\unsure{(}M),\\
 D_0\arrow[u,phantom,"\cup" description] \& {}
\end{tikzcd}}

associant à $\alpha$ le couple $(i,\unsure{j|p})$ (source et but de $\alpha$, considéré comme \unsure{flèche|flèches} dans $\underline{\mathrm{Fl}}(M)$ de $i$ dans $p$). La relation $R$ est alors déduite de $\varphi,D_0$ par la condition
% archive sheet number: 226
